\subsection{Gabarits}
\label{subsec:gabarits}

Pour s'assuré de la conformité de la fusée SP-02 aux exigences du \acrshort{CDC}, une analyse des gabarits de vol a été réalisée. Cette
analyse vise à vérifier que la fusée respecte les portées balistiques et les altitudes maximales autorisées pour chaque phase de vol, en
tenant compte des différentes configurations de vol possibles.\\

\begin{table}[h!]
    \centering
    \begin{tabular}{|l|l|l|l|}
        \hline
        \textbf{Etage} & \textbf{Etage 1} & \textbf{Etage 1 \& 2} & \textbf{Etage 2} \\
        \hline
        \textbf{Séparation} & Oui & Non & Oui \\
        \hline
        \textbf{Pro 54} & Actif & Actif & Actif \\
        \hline
        \textbf{Pro 24} & Inactif & Inactif & Actif \\
        \hline
        \textbf{Altitude z à 45° (m)} & 544.86 & 567.82 & 612.80 \\
        \hline
        \textbf{Portée x à 45° (m)} & 1932.77 & 2184.92 & 2153.82 \\
        \hline
        \textbf{Altitude z à 80° (m)} & 1228.51 & 1320.24 & 1439.27 \\
        \hline
        \textbf{Portée x à 80° (m)} & 737.71 & 859.25 & 839.71 \\
        \hline        
    \end{tabular}
    \caption{Résumé des gabarits de vol de la fusée SP-02 dans différentes configurations}
    \label{tab:gabarits_vol}
\end{table}